\begin{abstract}
Automatic Speech Recognition (ASR) foundation models achieve high accuracy on spoken prose but degrade catastrophically on singing voice. In tonal languages such as Vietnamese, this degradation is intensified by two coupled phenomena: (1) \textit{Melodic Tone Neutralization}, wherein musical melody ($F_0$) overrides lexical pitch contours, inducing severe tonal confusion; and (2) \textit{Attention Collapse}, where prolonged vibrato vowels trap autoregressive decoders in horizontal attention stagnation.

To resolve these challenges across the complete acoustic pipeline, we propose an integrated multi-tier framework: (i) front-end Glottal Inverse Filtering (IAIF) extracting source-filter parameters (Open Quotient $O_q$, spectral tilt $H_1^*-H_2^*$, and creaky phonation cues) that preserve lexical tone identity despite melody flattening; (ii) continuous pitch conditioning via gated semitone-cents embeddings coupled with Syllable-Tone InfoNCE contrastive regularization; (iii) a Multi-Dialect Mixture-of-Experts (MoE) architecture with dynamic gating across regional adapters; and (iv) a low-latency Streaming Monotonic Chunkwise Attention (MoChA) engine.

We validate our framework on a unified benchmark of 58 authentic audio files (434.65 seconds) across four domains: Vietnamese singing voice, Mozilla Common Voice 27.0, VietSuperSpeech conversational speech, and VIVOS studio speech. Empirical results demonstrate that our Multi-Dialect MoE architecture slashes singing Word Error Rate (WER) from $135.90\%$ to $54.12\%$ ($33.85\%$ under VietLyrics 647-hour macro pretraining), driving the global multi-domain Pareto score from $153.40\%$ down to $49.07\%$ (a $-104.33\%$ gain). 

Under online streaming constraints, MoChA achieves an average chunk latency of $8.50$\,ms on commodity laptop hardware (well below the $160$\,ms real-time budget) with a streaming Real-Time Factor (RTF) of $0.1091$ ($9.2\times$ faster than real-time) while retaining $98.72\%$ of offline recognition accuracy ($55.40\%$ singing WER).
\keywords{Singing Voice Recognition \and Automatic Lyric Transcription \and Tonal Languages \and Glottal Inverse Filtering \and Mixture-of-Experts \and Streaming MoChA.}
\end{abstract}
